\documentclass[11pt]{article}
\usepackage[a4paper,margin=26mm]{geometry}
\usepackage{amsmath,amssymb,mathtools}
\usepackage{unicode-math}
\setmainfont{Latin Modern Roman}
\setmathfont{Latin Modern Math}
\usepackage{xurl,hyperref}
\hypersetup{hidelinks,pdftitle={UnsolvedMath Problem 10400024}}
\newcommand{\sref}[2]{\hyperlink{source:#1:#2}{[S#2]}}
\newcommand{\cataloguescope}[1]{#1}
\begin{document}
% UnsolvedMath ID 10400024; source catalogue code AMR-103-0024
\setcounter{section}{10400023}
\section{Conjecture 2.2 — A(S1; Z) is torsion free.}\label{10400024}
{\small\sffamily UnsolvedMath ID 10400024 · Topology}
\subsection{Definitions and mathematical statement}

\paragraph{Shared notation.}$\mathbb N=\{1,2,\ldots\}$, and $\mathbb Z,\mathbb Q,\mathbb R,\mathbb C$ denote the integers, rationals, reals and complex numbers. Any other conventions are specified locally.


A Jacobi diagram on an oriented circle consists of a finite uni-trivalent graph with univalent vertices attached to the circle and a cyclic ordering of the three incident half-edges at each trivalent vertex. Diagrams are considered up to orientation-preserving identification of the circle and graph. Its degree is half the total number of univalent and trivalent vertices. Let $\mathcal A(S^1;R)$ be the graded $R$-module they generate modulo AS, IHX and STU: reversing a trivalent cyclic order changes sign (AS); the three local reconnections of four external half-edges satisfy the oriented Jacobi identity (IHX); the difference between the two orders of neighboring legs on the circle equals the diagram joining them at a trivalent vertex (STU). All relations keep the rest of the diagram fixed. The framing-independence relation FI additionally sets to zero a diagram with an isolated chord whose two endpoints are consecutive on the circle. Write $\mathcal A_d$ for degree $d$.
\paragraph{Statement recorded in the supplied source.}
Is $\mathcal A(S^1;\mathbb Z)$ torsion-free, equivalently is every degree piece torsion-free? Here no FI quotient is imposed: the question is for the integral AS--IHX--STU diagram module itself. Torsion-free means $m x=0$, $m\in\mathbb Z\setminus\{0\}$, implies $x=0$. \sref{10400024}{2}
\subsection{Short English statement}
Is the integral circle Jacobi-diagram module without framing-independence quotient free of torsion?
\subsection{Sources}
\begin{description}
\item[\hypertarget{source:10400024:1}{S1}] ULAM. AI and the UnsolvedMath Contributors, \emph{UnsolvedMath: A Curated Collection of Open Mathematics Problems}, record 10400024 (AMR-103-0024). CC BY 4.0. \url{https://huggingface.co/datasets/ulamai/UnsolvedMath}.
\item[\hypertarget{source:10400024:2}{S2}] T. Ohtsuki (editor), \emph{Problems on invariants of knots and 3-manifolds}, Geometry \& Topology Monographs 4 (2002), pp. 377--572 (published 2004), item 2.2 and its surrounding definitions and remarks. \url{https://arxiv.org/abs/math/0406190v1}.
\end{description}
\label{end:10400024}
\end{document}
